\BourbakiExerciseGroup

All rings under consideration are assumed to be commutative unless explicitly stated otherwise.

\begin{enumerate}
\def\labelenumi{\arabic{enumi})}
\tightlist
\item
  Let A be a totally ordered ring and B an additive subgroup of A which is closed under multiplication.
\end{enumerate}

\begin{enumerate}
\def\labelenumi{\alph{enumi})}
\item
  An element \(x \in A\) is said to be infinitely large with respect to B if \(|y| < |x|\) for all \(y \in B\) . Show that the set \(F(B)\) of elements of A which are not infinitely large with respect to B is a subring of A containing B.
\item
  An element \(x \in A\) is said to be infinitely small with respect to B if, for all y \textgreater{} 0 in B, we have \(|x| < y\) . If, for all y \textgreater{} 0 in B, there exists \(z \in B\) such that 0 \textless{} z \textless{} y, then the set of elements of A which are infinitely small with respect to B is a subpseudo-ring I(B) of A. If in addition, for every pair of elements y, z of B with 0 \textless{} y \textless{} z, there exists \(x \in B\) with 0 \textless{} xz \textless{} y, then I(B) is an ideal in the ring F(B).
\end{enumerate}

\begin{enumerate}
\def\labelenumi{\arabic{enumi})}
\setcounter{enumi}{1}
\item
  Let A be a totally ordered ring and n the set of nilpotent elements of A (which is an ideal of A). Every element of A not belonging to n is infinitely large with respect to n; the quotient ring A/n is totally ordered (VI, p.~30, Ex. 4) and is an integral domain.
\item
  In the field \(\mathbf{Q}\) , let \(\mathbf{P}\) be the set consisting of 0 and the rational numbers \(\geqslant 1\) (in the usual ordering). Show that \(\mathbf{P}\) satisfies axioms \((\mathrm{AP}_{\mathrm{I}})\) , \((\mathrm{AP}_{\mathrm{II}})\) and \((\mathrm{AP}_{\mathrm{III}})\) .
\item
  \begin{enumerate}
  \def\labelenumii{\alph{enumii})}
  \tightlist
  \item
    Let K be a field, and P a subset of K satisfying conditions \((\mathrm{AP}_{\mathrm{I}})\) , \((\mathrm{AP}_{\mathrm{II}})\) and \((\mathrm{AP}_{\mathrm{III}})\) , and such that \(K^{2} \subset P\) (where \(K^{2}\) denotes the set of squares of elements of K); show that if x \textgreater{} 0 in the ordering defined by P, then \(x^{-1} > 0\) ; deduce that the set \(K_{+}^{*}\) of elements \textgreater{} 0 in K is a subgroup of the multiplicative group of K.
  \end{enumerate}
\end{enumerate}

\begin{enumerate}
\def\labelenumi{\alph{enumi})}
\setcounter{enumi}{1}
\item
  In the field \(\mathbf{Q}\) , the only set \(\mathbf{P}\) which satisfies the conditions of \(a\) ) is the set of rational numbers \(\geqslant 0\) (in the usual ordering).
\item
  * Let P be a subset of K satisfying \((\mathrm{AP}_{\mathrm{I}})\) , \((\mathrm{AP}_{\mathrm{II}})\) and \((\mathrm{AP}_{\mathrm{III}})\) , such that \(1 \in P\) and such that x \textgreater{} 0 implies \(x^{-1} > 0\) in the ordering defined by P. Show that \(y^{2} \geqslant z^{2}\) implies \(y \geqslant z\) , for positive y and z. Deduce that, for arbitrary y \textgreater{} 0, we have
\end{enumerate}

\[
\frac {1}{2} (y + y ^ {- 1}) \geqslant 1 - n ^ {- 1}
\]

for every integer \(n > 0\) (notice that \((y + y^{-1})^{2m} \geqslant \binom{2m}{m}\) for every integer \(m > 0\) ). Deduce that, if the ordered additive group structure defined by P is archimedean (VI, p.~35, Ex. 31), then \((y - z)^2 \geqslant 0\) for every pair of elements \(y, z\) of P; if \(K'\) is the subfield of K generated by P then \(x^2 \geqslant 0\) for all \(x \in K'\) . *

\(d) * \text{Let } K = R(X) \text{ be the field of rational functions in one indeterminate over } R; \text{ let } P \text{ be the set consisting of 0 and the rational functions } u \in K \text{ such that } u(t) \text{ is defined and } >0 \text{ for every real number } t. \text{ Show that } P \text{ satisfies the conditions of } c) \text{ and generates } K, \text{ but that there exist elements } v \in K \text{ such that } v^2 \notin P. *\)

\begin{enumerate}
\def\labelenumi{\arabic{enumi})}
\setcounter{enumi}{4}
\tightlist
\item
  Let A be a lattice ordered ring. A convex ideal (VI, p.~30, Ex. 4) of A is called irreducible if it is not the intersection of two convex ideals distinct from itself.
\end{enumerate}

\begin{enumerate}
\def\labelenumi{\alph{enumi})}
\tightlist
\item
  Show that the intersection of the irreducible convex ideals of A is 0 (if \(a \in \mathbf{A}\) is nonzero, consider a convex ideal of A maximal among those which do not contain \(a\) ). Deduce that A is isomorphic to a subring \(\mathbf{A}'\) of a product \(\prod_{i} \mathbf{A}_i\) , such that each \(\mathbf{A}_i\) is lattice ordered,
\end{enumerate}

\(\mathrm{pr}_{\iota}(\mathbf{A}^{\prime})=\mathbf{A}_{\iota}\) for all \(\iota\) and the ideal \(\{0\}\) is irreducible in \(\mathbf{A}_{\iota}\) (as a convex ideal) for all \(\iota\) . b) Show that the following conditions are equivalent in A:

α) \(|xy| = |x| \cdot |y|\) for all x, y;

\(\beta) x \cdot \sup(y, z) = \sup(xy, xz)\) for all \(x \geqslant 0\) , \(y, z\) ;

\(\gamma) x \cdot \inf(y, z) = \inf(xy, xz)\) for all \(x \geqslant 0\) , \(y, z\) ;

\(\delta)\inf (y,z) = 0\) implies \(\inf (xy,z) = 0\) for all \(x\geqslant 0\) .

(To see that \(\gamma\) ) implies \(\delta\) ), observe that \(xy \leqslant y \cdot \sup (x, 1)\) ).

When these conditions are satisfied, then A is said to be strongly lattice ordered.

\begin{enumerate}
\def\labelenumi{\alph{enumi})}
\setcounter{enumi}{2}
\tightlist
\item
  Show that a ring \(\mathbf{A}\) is strongly lattice ordered if and only if it is isomorphic to a subring of a product \(\prod_{i} \mathbf{A}_{i}\) of totally ordered rings. (Using \(a\) ), reduce to showing that if
\end{enumerate}

\{0\} is irreducible in a strongly lattice ordered ring A, then A is totally ordered; for this, note that if B is a strongly lattice ordered ring then for all \(b \in B\) the set \(m\) of \(x \in B\) such that \(\inf(|x|, |b|) = 0\) and the set \(n\) of \(y \in B\) such that \(|y| \leqslant |zb|\) for some \(z \in B\) are two convex ideals of B such that \(m \cap n = \{0\}\) .

\(d\) ) Show that in a strongly lattice ordered ring \(\mathbf{A}\) the relation \(\inf (x,y) = 0\) implies \(xy = 0\) ; for all \(z\in \mathbf{A}\) we have \(z^2\geqslant 0\) ; for all \(x,y\in \mathbf{A}\) we have \(xy\leqslant \sup (x^2,y^2)\)

\begin{enumerate}
\def\labelenumi{\alph{enumi})}
\setcounter{enumi}{4}
\tightlist
\item
  Let A be a lattice ordered ring with no nilpotent element \textgreater0; show that if \(\inf(x, y) = 0\) implies xy = 0 then A is strongly lattice ordered (check condition \(\delta\) ) of b)).
\end{enumerate}

\begin{enumerate}
\def\labelenumi{\arabic{enumi})}
\setcounter{enumi}{5}
\tightlist
\item
  \begin{enumerate}
  \def\labelenumii{\alph{enumii})}
  \tightlist
  \item
    Let K be a lattice ordered field. Show that the following conditions are equivalent:
  \end{enumerate}
\end{enumerate}

\(\alpha) x^{2} \geqslant 0\) for all \(x \in \mathbf{K}\) ;

\(\beta) x > 0\) implies \(x^{-1} > 0\) ;

γ) K is strongly lattice ordered (Ex. 5);

δ) K is totally ordered.

(To see that \(\beta\) ) implies \(\underline{\delta}\) ), notice that \(\beta\) implies \(xy \leqslant x^2 + y^2\) for \(x, y \in K\) ).

\begin{enumerate}
\def\labelenumi{\alph{enumi})}
\setcounter{enumi}{1}
\item
  Let K be the field \(\mathbf{Q}(\sqrt{2})\) obtained by adjoining \(\sqrt{2}\) to the field of rational numbers; as an additive group K is identified with \(Q \times Q\) via the bijection \(x + y \sqrt{2} \mapsto (x, y)\) ; show that if the product ordering on \(Q \times Q\) (where Q has the usual ordering) is carried over to K via this bijection, then K is lattice ordered but not strongly lattice ordered.
\item
  The field \(\mathbf{K} = \mathbf{Q}(\mathbf{X})\) of rational functions in one indeterminate X over Q is generated by its set of squares \(K^{2}\) . Show that the set P of sums of squares of elements of K defines a non-lattice ordering on K compatible with its ring structure, and such that u \textgreater{} 0 implies \(u^{-1} > 0\) .
\end{enumerate}

\begin{enumerate}
\def\labelenumi{\arabic{enumi})}
\setcounter{enumi}{6}
\item
  Every element of a commutative field K of characteristic \(\neq 2\) is a sum of squares if and only if -1 is a sum of squares (notice that every element of K is a difference of two squares).
\item
  Let A be a nonempty subset of a commutative field K of characteristic \(\neq 2\) ; then K is said to be A-orderable if no element of A is a sum of squares in K. We say that K is orderable if there exists a total ordering of K which is compatible with its ring structure.
\end{enumerate}

\begin{enumerate}
\def\labelenumi{\alph{enumi})}
\item
  Show that if K is A-orderable then it is orderable (Ex. 7), and hence of characteristic 0.
\item
  Show that pure extensions and algebraic extensions of odd degree of an A-orderable field K are A-orderable (argue as in Prop. 3 of VI, p.~23).
\item
  Let K be an A-orderable field and let b be an element of K not of the form ca - d, where \(a \in A\) and c and d are sums of squares in K. Show that the field \(\mathrm{K}(\sqrt{b})\) is A-orderable.
\item
  An A-orderable field K is called maximal if no proper algebraic extension of K is A-orderable. Show that a maximal A-orderable field K has the following properties: 1. K is pythagorean, in other words every sum of squares is a square; 2. no element of A is a square; 3. every element of K which is not a square has the form \(c^2 a - d^2\) , where \(a \in \mathbf{A}\) ; 4. every odd degree polynomial in K{[}X{]} has at least one root in K (use b) and c)). e) Show that if K satisfies the conditions of d) then every element algebraic over K has degree \(2^q\) over K for some q (argue as in Prop. 8 of VI, p.~26, by induction on the exponent of 2 in the degree of the element under consideration). Show that on the other hand no quadratic extension of K is A-orderable. Deduce that K is a maximal A-orderable field (use Galois theory and Prop. 12 of I, p.~77).
\item
  Let K be an A-orderable field and \(\Omega\) an algebraically closed extension of K. Show that there exists a maximal A-orderable field contained in \(\Omega\) and containing K.
\end{enumerate}

\begin{enumerate}
\def\labelenumi{\arabic{enumi})}
\setcounter{enumi}{8}
\tightlist
\item
  Let \(q > 0\) be a non-square natural number; let \(A\) denote the set \(\{-1, \sqrt{q}\}\) in the field \(K = Q(\sqrt{q})\) ; show that \(K\) is \(A\) -orderable. Show that there exists an extension \(E\) of \(K\) such that \(E\) is pythagorean, orderable, and that every odd degree polynomial over \(E\) has a root in \(E\) , but that \(E\) does not admit the structure of a maximal ordered field (consider a maximal \(A\) -orderable extension of \(K\) ).
\end{enumerate}

¶ 10) a) Let K be an orderable field and E a Galois extension of K. Show that either E is orderable or there exists an orderable algebraic extension F of K, contained in E, such that E is a quadratic extension of F (use Th. 6 of V, p.~70).

\begin{enumerate}
\def\labelenumi{\alph{enumi})}
\setcounter{enumi}{1}
\tightlist
\item
  Show that the polynomial \(X^{4} + 2\) is irreducible over Q, and that the extension K of Q of degree 4, obtained by adjoining a root of this polynomial to Q, contains no orderable subfield other than Q (find all the subfields of K by Galois theory).
\end{enumerate}

\begin{enumerate}
\def\labelenumi{\arabic{enumi})}
\setcounter{enumi}{10}
\tightlist
\item
  Let K be an ordered field and G a subfield of K. a) Show that \(x \neq 0\) is infinitely large with respect to G if and only if \(x^{-1}\) is infinitely small with respect to G. We say that K is comparable to G if there is no element of K which is infinitely large with respect to G (nor, consequently, any nonzero element of K which is infinitely small with respect to G). A necessary and sufficient condition for K to be comparable to its prime subfield Q is that K be archimedean (VI, p.~35, Ex. 31), * and hence that K be isomorphic to a subfield of R (VI, p.~36, Ex. 33c)).
\end{enumerate}

\begin{enumerate}
\def\labelenumi{\alph{enumi})}
\setcounter{enumi}{1}
\item
  Show that, in the ring \(F(G)\) of elements of K which are not infinitely large with respect to G, the set \(I(G)\) of elements infinitely small with respect to G is a maximal ideal; moreover the ordering on the quotient field \(K(G) = F(G)/I(G)\) induced from that of \(F(G)\) is compatible with the ring structure and is total.
\item
  Show that a class modulo \(I(G)\) can contain at most one element of G; deduce that the natural map from G to \(K(G)\) is an isomorphism from the ordered field G onto a subfield \(G'\) of \(K(G)\) , and that \(K(G)\) is comparable to \(G'\) .
\end{enumerate}

\begin{enumerate}
\def\labelenumi{\arabic{enumi})}
\setcounter{enumi}{11}
\item
  Let K be a maximal ordered field, let f be a polynomial over K and let a and b be two roots of f such that a \textless{} b and such that f has no roots between a and b. Show that if g is a rational function over K whose denominator does not vanish for \(a \leqslant x \leqslant b\) then the equation \(f(x)g(x) + f'(x) = 0\) has an odd number of solutions in \([a, b]\) (where each solution is counted with its multiplicity; use Prop. 5 of VI, p.~25). Deduce that if h is a rational function over K having a and b as roots and whose denominator does not vanish in \([a, b]\) , then the equation \(h'(x) = 0\) has at least one root in \([a, b]\) .
\item
  Let K be a maximal ordered field, let h be a rational function over K and let \([a, b]\) be a closed interval in which h is defined. Show that there exists \(c \in ]a, b[\) such that \(h(b) - h(a) = (b - a)h'(c)\) (use Ex. 12). Deduce that h is an increasing function in \([a, b]\) if and only if \(h'(x) \geqslant 0\) in this interval (to see that this condition is necessary, split the interval by the roots of \(h'(x) = 0\) ).
\item
  \begin{enumerate}
  \def\labelenumii{\alph{enumii})}
  \tightlist
  \item
    Let K be a maximal ordered field, let E be a subfield of K and let f be a polynomial over E; show that all the roots of f in K belong to F(E) (Ex. 11, b)) (use Prop. 4 of VI, p.~24). Deduce that if G is a subfield of K and E ⊂ K is an extension comparable to G (Ex. 11) then the set of elements of K algebraic over E is a maximal ordered field comparable to G.
  \end{enumerate}
\end{enumerate}

\begin{enumerate}
\def\labelenumi{\alph{enumi})}
\setcounter{enumi}{1}
\item
  Deduce from a) that, under the hypotheses of a), the field K(G) (Ex. 11, b)) is a maximal ordered field.
\item
  Let \(f\) be a polynomial of degree \(\geqslant 1\) over \(G\) . Show that \(f(t)\) is infinitely large with respect to \(G\) if and only if \(t\) is infinitely large with respect to \(G\) ; the element \(f(t)\) is infinitely small with respect to \(G\) if and only if \(t\) is congruent modulo \(I(G)\) to a root of \(f\) in \(K\) (split \(K\) into intervals using the roots of \(f\) and \(f'\) in \(K\) , and apply Ex. 13; notice that if \(x < t\) and if \(x \in K\) is not congruent modulo \(I(G)\) to \(t\) , then there exists \(y \in K\) such that \(x < y < t\) and \(y - x \in G\) ).
\end{enumerate}

\begin{enumerate}
\def\labelenumi{\arabic{enumi})}
\setcounter{enumi}{14}
\tightlist
\item
  Let E be an ordered field, let K be a subfield of E such that E is algebraic over K, and let R be a maximal ordered extension of K.
\end{enumerate}

\begin{enumerate}
\def\labelenumi{\alph{enumi})}
\item
  Among the elements of E - K, let \(x_0\) be one whose minimal polynomial \(f\) over K has the least possible degree. Show that \(f'(T)\) is a product of factors in R{[}T{]} of the form \((T - a_i)^2 + c_i^2\) ( \(a_i, c_i \in \mathbb{R}\) ), and first order factors \(T - b_j\) with the \(b_j\) pairwise distinct elements of K. Deduce that there exists an element \(y_0\) in R such that \(f(y_0) = 0\) and such that \(x_0 - z\) and \(y_0 - z\) have the same sign, in E and R respectively, for all \(z \in \mathbb{K}\) (cf.~VI, p.~43, Ex. 26, c)).
\item
  Deduce from a) that there exists a K-isomorphism from the ordered field E onto a subfield of R. (First show that there exists such an isomorphism for the subfield \(\mathrm{K}(x_{0})\) of E; for any polynomial \(g \in \mathrm{K}[\mathrm{T}]\) such that \(\deg(g) < \deg(f)\) , apply the same argument used for \(f'\) in a) to evaluate the signs of \(g(x_{0})\) and \(g(y_{0})\) ; then apply Zorn's Lemma).
\item
  Show that if R and R' are two maximal ordered algebraic extensions of K then there exists precisely one K-isomorphism from R onto R' (apply b)).
\end{enumerate}

\begin{enumerate}
\def\labelenumi{\arabic{enumi})}
\setcounter{enumi}{15}
\tightlist
\item
  Let K be a maximal ordered field and let
\end{enumerate}

\[
g (\mathrm{T}) = a _ {0} + a _ {1} \mathrm{T} + \dots + a _ {n} \mathrm{T} ^ {n}
\]

be a nonzero polynomial in K{[}T{]} whose roots all lie in K. Show that for any polynomial \(f \in K[T]\) , the number of roots of the polynomial

\[
a _ {0} f (\mathrm{T}) + a _ {1} f ^ {\prime} (\mathrm{T}) + a _ {2} f ^ {\prime \prime} (\mathrm{T}) + \dots + a _ {n} f ^ {(n)} (\mathrm{T})
\]

which do not belong to K, counted with multiplicity, is at most equal to the number of roots of f, counted with multiplicity, which do not belong to K (use Ex. 12 for n = 1, then proceed by induction). Deduce that the polynomial

\[
a _ {0} + \frac {a _ {1}}{1 !} \mathrm{T} + \frac {a _ {2}}{2 !} \mathrm{T} ^ {2} + \dots + \frac {a _ {n}}{n !} \mathrm{T} ^ {n}
\]

has all its roots in K.

\begin{enumerate}
\def\labelenumi{\arabic{enumi})}
\setcounter{enumi}{16}
\tightlist
\item
  Let K be a maximal ordered field, let g be a nonzero polynomial in K{[}T{]} whose roots are all in K, and do not belong to the interval \([0, n]\) of K. Show that for every polynomial \(f(T) = a_{0} + a_{1}T + \cdots + a_{n}T^{n}\) of K{[}T{]}, the number of roots of the polynomial
\end{enumerate}

\[
a _ {0} g (0) + a _ {1} g (1) \mathrm{T} + a _ {2} g (2) \mathrm{T} ^ {2} + \dots a _ {n} g (n) \mathrm{T} ^ {n}
\]

which do not belong to K, counted with multiplicity, is at most equal to the number of roots of f, counted with multiplicity, which do not belong to K (same method as Ex. 16).

\begin{enumerate}
\def\labelenumi{\arabic{enumi})}
\setcounter{enumi}{17}
\item
  Let K be a maximal ordered field and let f be a polynomial of degree n in K{[}T{]}, all of whose roots are in K. Show that for all \(c \neq 0\) in K the polynomial \(f^{2} + cf'\) has at least n - 1 and at most \(n + 1\) roots in K, counted with multiplicity (use Ex. 13).
\item
  Let K be a maximal ordered field and \(f\) a nonzero polynomial in K{[}T{]}. A necessary and sufficient condition that, for every polynomial \(g \in \mathbf{K}[\mathbf{T}]\) , the number of roots of \(fg + g'\) , counted with multiplicity, not belonging to K, be at most equal to the number of roots of \(g\) , counted with multiplicity, not belonging to K, is that \(f(\mathbf{T}) = a - b\mathbf{T}\) with \(b \geqslant 0\) in K. (To see that the condition is sufficient use Ex. 12; to show that it is necessary, take g = 1 and g = f).
\item
  Let \(\mathbf{K}\) be a maximal ordered field and
\end{enumerate}

\[
f (\mathrm{T}) = a _ {0} + \binom {n} {1} a _ {1} \mathrm{T} + \binom {n} {2} a _ {2} \mathrm{T} ^ {2} + \dots + a _ {n} \mathrm{T} ^ {n}
\]

a polynomial in K{[}T{]}. For all integers p, q such that \(0 \leqslant p < p + q \leqslant n\) , the number of roots of the polynomial

\[
a _ {p} + \binom{q}{1} a _ {p + 1} \mathrm{T} + \binom{q}{2} a _ {p + 2} \mathrm{T} ^ {2} + \dots + a _ {p + q} \mathrm{T} ^ {q}
\]

counted with multiplicity, which do not belong to K is at most equal to the number of roots of f, counted with multiplicity, which do not belong to K (use Ex. 12).

Deduce that if \(b_{1}, \ldots, b_{n}\) are n distinct elements \textgreater0 in K, and if

\[
\left(\mathrm{T} + b _ {1}\right) \dots \left(\mathrm{T} + b _ {n}\right) = \mathrm{T} ^ {n} + \binom {n} {1} m _ {1} \mathrm{T} ^ {n - 1} + \dots + m _ {n}
\]

then

\[
m _ {k} ^ {1 / k} > m _ {k + 1} ^ {1 / (k + 1)} \quad \text { for } \quad 1 \leqslant k \leqslant n - 1.
\]

\begin{enumerate}
\def\labelenumi{\arabic{enumi})}
\setcounter{enumi}{20}
\tightlist
\item
  Let \((a_{i})_{1 \leqslant i \leqslant n}\) be a finite sequence of n elements of an ordered field K, not all zero; let \((a_{i_{k}})_{1 \leqslant k \leqslant p}\) be the subsequence of \((a_{i})\) consisting of the nonzero \(a_{i}\) ( \(i_{1} < i_{2} < \cdots < i_{p}\) ); then the number of indices \(k \leqslant p - 1\) such that \(a_{i_{k}}\) and \(a_{i_{k+1}}\) have opposite signs is called the number of variations of the sequence \((a_{i})\) .
\end{enumerate}

Let K be a maximal ordered field and let f be a nonzero polynomial in K{[}T{]} of degree n \textgreater{} 0; for all \(a \in K\) let \(w(a)\) denote the number of variations of the sequence \((f^{(i)}(a))_{0 \leq i \leq n}\) . If v is the number of roots of f, counted with multiplicity, in the interval \([a, b]\) (a \textless{} b), show that \(v \leq w(a) - w(b)\) and that \(w(a) - w(b) - v\) is even (« Budan-Fourier rule »; split the interval \([a, b]\) by the roots of f, and evaluate the amount of variation in \(w(x)\) as x passes one of these roots).

Deduce that if \(f(T) = a_{0} + a_{1}T + \cdots + a_{n}T^{n}\) , then the number of roots \textgreater{} 0 of f, counted with multiplicity, is at most equal to the number of variations of the sequence \((a_{i})_{0 \leq i \leq n}\) , and that the difference between these two numbers is even (« Descartes' rule »).

\begin{enumerate}
\def\labelenumi{\arabic{enumi})}
\setcounter{enumi}{21}
\tightlist
\item
  Let K be a maximal ordered field and let
\end{enumerate}

\[
f (\mathrm{T}) = a _ {0} + a _ {1} \mathrm{T} + \dots + a _ {n} \mathrm{T} ^ {n}
\]

be a polynomial in K{[}T{]} such that \(a_0 \neq 0 \neq a_n\) and

\[
a _ {p} = a _ {p + 1} = \dots = a _ {p + 2 m - 1} = 0 \quad (1 \leqslant p <   p + 2 m - 1 <   n).
\]

Show that \(f\) has at most \(n - 2m\) roots in \(\mathbf{K}\) , counted with multiplicity (apply Descartes' rule).

Deduce that if \(g(T) = 1 + c_{1}T + \cdots + c_{n}T^{n}\) has all its roots in K, and if \(h(T) = 1 + b_{1}T + \cdots + b_{2m}T^{2m}\) is the polynomial consisting of the first \(2m + 1\) terms of the formal power series \(1/g \in K[[T]]\) , then h has no roots in K.

\begin{enumerate}
\def\labelenumi{\arabic{enumi})}
\setcounter{enumi}{22}
\tightlist
\item
  Let K be a maximal ordered field and let \(a_{1}, \ldots, a_{n}, b_{1}, \ldots, b_{n}\) be 2n elements of K such that \(\sum a_{i} \neq 0\) and \(b_{1} < b_{2} < \cdots < b_{n}\) . Show that for any integer \(m \geqslant 1\) the number v of roots in K, counted with multiplicity, of the polynomial
\end{enumerate}

\[
f (\mathrm{T}) = a _ {1} \left(\mathrm{T} - b _ {1}\right) ^ {m} + a _ {2} \left(\mathrm{T} - b _ {2}\right) ^ {m} + \dots + a _ {n} \left(\mathrm{T} - b _ {n}\right) ^ {m}
\]

is at most equal to the number of variations w of the sequence

\[
\left(a _ {1}, a _ {2}, \dots , a _ {n}, (- 1) ^ {m} a _ {1}\right),
\]

and that the difference w - v is even. (Argue by induction on n, applying Ex. 13 (VI, p.~40) to a rational function of the form \(f(\mathrm{T})/(\mathrm{T}-c)^{m}\) .)

\begin{enumerate}
\def\labelenumi{\arabic{enumi})}
\setcounter{enumi}{23}
\item
  Let K be a maximal ordered field; consider an ordering on the field K(X) of rational functions which makes it an ordered extension of K. Show that this ordering is determined by the set A of \(x \in K\) such that x \textless{} X (show that the sign of every polynomial f over K is determined by A, using Prop. 9 of VI, p.~28). Conversely, show that to every set A ⊂ K such that \(y \in A\) whenever \(y \leqslant x\) and \(x \in A\) , there corresponds an ordering of K(X) making it an ordered extension of K and such that A is the set of \(x \in K\) such that x \textless{} X (same method). When A has a greatest element, or CA has a least element, or when A = K or A = ∅, then the corresponding ordering of K(X) is such that K(X) is not comparable to K. On the other hand, if neither A nor CA is empty, and there is no greatest element in A nor least element in CA, then K(X) is comparable to K.
\item
  Use Ex. 24 to find an example of a field E admitting several ordered field structures, none of which is induced from another via an automorphism of E. Hence find an example of a field, equipped with an ordering compatible with its ring structure but not a lattice ordering (consider the diagonal in E × E and use VI, p.~38, Ex. 5).
\item
  \begin{enumerate}
  \def\labelenumii{\alph{enumii})}
  \tightlist
  \item
    If K is an archimedean ordered field (VI, p.~40, Ex. 11, a)), and x and y are two elements of K with x \textless{} y, show that there exists \(r \in Q\) such that x \textless{} r \textless{} y. * Deduce that there is only one ordered subfield of R isomorphic to K. *
  \end{enumerate}
\end{enumerate}

\begin{enumerate}
\def\labelenumi{\alph{enumi})}
\setcounter{enumi}{1}
\tightlist
\item
  Consider the ordering of the field \(\mathbf{K} = \mathbf{Q}(\mathbf{X})\) in which \(\mathbf{X}\) is infinitely large with respect to \(\mathbf{Q}\) (Ex. 24). Show that \(\mathbf{K}\) is comparable to the subfield \(\mathbf{Q}(\mathbf{X}^2)\) , and give an example of two elements \(x, y\) of \(\mathbf{K}\) such that \(x < y\) and no element of \(\mathbf{Q}(\mathbf{X}^2)\) lies between \(x\) and \(y\) . c) Let \(\mathbf{K}\) be the ordered field defined in \(b\) ; show that the polynomial
\end{enumerate}

\[
f (\mathbf {Y}) = (\mathbf {Y} ^ {2} - \mathbf {X}) (\mathbf {Y} ^ {2} - 4 \mathbf {X}) - 1
\]

over K is irreducible, that it has roots in every maximal ordered extension of K, and that \(f(a) > 0\) for all \(a \in K\) .

\begin{enumerate}
\def\labelenumi{\arabic{enumi})}
\setcounter{enumi}{26}
\tightlist
\item
  Let K be an ordered field, let E be a purely transcendental extension of K and let \((\mathbf{X}_{i})_{i\in I}\) be a transcendence basis for E (V, p.~109).
\end{enumerate}

\begin{enumerate}
\def\labelenumi{\alph{enumi})}
\item
  * If K is archimedean then E has a structure of an ordered extension of K for which E is comparable to K if and only if the cardinality of I is at most equal to that of a transcendence basis for R over K (where K is considered as an ordered subfield of R, cf.~Ex. 26, a)); the set of such structures is then bijective with the set of injective maps from I into R such that \(f(I)\) is algebraically independent over K. *
\item
  If K is not archimedean then there always exists (at least) one structure on E of an ordered extension of K such that E is comparable to K (when I contains only one element, use Ex. 24 and notice that Q does not have a supremum in K; pass to the general case by Zorn's Lemma).
\end{enumerate}

* 28) a) Let K be a subfield of R and let \(\theta\) be a real number, algebraic over K. Show that the number of structures on K(θ) of an ordered extension of K is equal to the number of real conjugates of \(\theta\) (use Ex. 14, a) (VI, p.~40) and 26, a)).

\begin{enumerate}
\def\labelenumi{\alph{enumi})}
\setcounter{enumi}{1}
\tightlist
\item
  Let K be a subfield of R, equipped with n distinct structures of an ordered extension of Q, all archimedean; let \(P_{i}\) ( \(1 \leqslant i \leqslant n\) ) be the sets of positive elements under these orderings. Show that there exists a family of n elements \(b_{i}\) of K ( \(1 \leqslant i \leqslant n\) ) such that \(b_{i} \in P_{i}\) for all i and \(b_{i} \notin P_{k}\) for all \(k \neq i\) . (For each pair of distinct indices i, k, let \(a_{ik} \in P_{k}\) be such that \(a_{ik} \notin P_{i}\) ; consider the elements
\end{enumerate}

\[
c _ {i} = \sum_ {k \neq i} \left(\frac {1 + a _ {i k}}{1 - a _ {i k}}\right) ^ {2 r}
\]

where \(r\) is a sufficiently large integer.)

\begin{enumerate}
\def\labelenumi{\alph{enumi})}
\setcounter{enumi}{2}
\tightlist
\item
  Let \(\mathbf{K} = \mathbf{Q}(\theta)\) be an algebraic extension of \(\mathbf{Q}\) contained in \(\mathbf{R}\) , and let \(n\) be the number of real conjugates of \(\theta\) . Let \(\mathbf{C}\) be the set of sums of squares of elements of \(\mathbf{K}\) ; show that \(\mathbf{C}^* = \mathbf{C} \cap \mathbf{K}^*\) is a subgroup of \(\mathbf{K}^*\) and that \((\mathbf{K}^*: \mathbf{C}^*) = 2^n\) . (Use \(b\) ) and Cor. 2 of VI, p.~23.) *
\end{enumerate}

\begin{enumerate}
\def\labelenumi{\arabic{enumi})}
\setcounter{enumi}{28}
\item
  Let K be a maximal ordered field and let G be a subfield of K. Show that the set of extensions of G contained in K, and which are comparable to G, is inductively ordered; if \(E_{0}\) is a maximal element of this set, show that \(E_{0}\) is isomorphic to the field \(\mathrm{K}(\mathrm{G})\) defined in Ex. 11, b) of VI, p.~40 (prove that the natural map from \(\mathrm{F}(\mathrm{G})\) onto \(\mathrm{K}(\mathrm{G})\) maps \(E_{0}\) onto \(\mathrm{K}(\mathrm{G})\) , by first showing that \(E_{0}\) is a maximal ordered field, using Ex. 14, a) of VI, p.~40, and then that no element of \(\mathrm{K}(\mathrm{G})\) is transcendental over the image of \(E_{0}\) , using Ex. 24).
\item
  \begin{enumerate}
  \def\labelenumii{\alph{enumii})}
  \tightlist
  \item
    Let K be a maximal ordered field and let m and M be two elements of K with m \textless{} M. Show that every polynomial \(f \in K[X]\) which is positive in the interval \((m, M)\) is a sum of polynomials of the form \((\alpha X + \beta) g^{2}\) , where \(g \in K[X]\) and \(\alpha X + \beta\) is positive in \((m, M)\) . (Reduce to the case of first or second degree polynomials, and use the formulae
  \end{enumerate}
\end{enumerate}

\[
\begin{array}{l} (\mathbf {X} - a) (\mathbf {X} - b) = (\mathbf {X} - b) ^ {2} + (b - a) (\mathbf {X} - b) \\ (\mathbf {X} - a) (b - \mathbf {X}) = ((\mathbf {X} - a) (b - \mathbf {X}) ^ {2} + (b - \mathbf {X}) (\mathbf {X} - a) ^ {2}) / (b - a) \end{array}
\]

for \(a < b\) .)

\begin{enumerate}
\def\labelenumi{\alph{enumi})}
\setcounter{enumi}{1}
\tightlist
\item
  Show that the result of a) is no longer necessarily true when K is an arbitrary ordered field (notice that a polynomial can be \textgreater0 in K, but not in a maximal ordered extension of K; cf.~Ex. 26, c)).
\end{enumerate}

\begin{enumerate}
\def\labelenumi{\arabic{enumi})}
\setcounter{enumi}{30}
\tightlist
\item
  \begin{enumerate}
  \def\labelenumii{\alph{enumii})}
  \tightlist
  \item
    Let \(\mathbf{K}\) be a field whose algebraic closure \(\mathbf{E}\) is an extension of \(\mathbf{K}\) of prime degree \(q\) . Show that \(\mathbf{K}\) is perfect (V, p.~7).
  \end{enumerate}
\end{enumerate}

\begin{enumerate}
\def\labelenumi{\alph{enumi})}
\setcounter{enumi}{1}
\item
  Show that q is distinct from the characteristic of K (V, p.~162, Ex. 9).
\item
  Show that K contains the q-th roots of unity, and that E is the splitting field of an irreducible polynomial of the form \(X^{q}-a\) over K; deduce that q=2 (otherwise, use Ex. 5, V, p.~161 to deduce that \(X^{q^{2}}-a\) would be irreducible). Show also that -a is a square in K (V, loc. cit.), that -1 is not a square in K, and that \(E = K(i)\) ( \(i^{2} = -1\) ).
\item
  Now suppose that K is such that its algebraic closure E has arbitrary finite degree \(\neq1\) over K. Show that \(i\notin K\) and that \(\mathrm{E}=\mathrm{K}(i)\) (if \(\mathrm{E}\neq\mathrm{K}(i)\) , show by Galois theory that there would exist a field F such that \(\mathrm{K}(i)\subset\mathrm{F}\subset\mathrm{E}\) , and that E has prime degree over F; then apply c)). Deduce that K is a maximal ordered field (under some suitable ordering) (show by induction that any sum of n squares in K is a square in K; to prove that \(a^{2}+b^{2}\) is a square, consider a square root \(x+iy\) of \(a+ib\) in \(\mathrm{K}(i)\) ).
\end{enumerate}

\begin{enumerate}
\def\labelenumi{\arabic{enumi})}
\setcounter{enumi}{31}
\tightlist
\item
  Let A be an algebraically closed field of characteristic 0.
\end{enumerate}

\begin{enumerate}
\def\labelenumi{\alph{enumi})}
\item
  The only elements of finite order in the group \(\operatorname{Aut}(\mathbf{A})\) of automorphisms of \(\mathbf{A}\) are the identity element and elements of order 2 (involutions of \(\mathbf{A}\) ); these latter correspond bijectively to the maximal ordered subfields \(\mathbf{E}\) of \(\mathbf{A}\) such that \([\mathbf{A}:\mathbf{E}] = 2\) , that is \(\mathbf{A} = \mathbf{E}(i)\) (cf.~Ex. 31). If \(\sigma\) is an involution of \(\mathbf{A}\) and \(\mathbf{E}\) is the fixed subfield of \(\sigma\) , then the group \(\operatorname{Aut}(\mathbf{E})\) of automorphisms of \(\mathbf{E}\) is isomorphic to the quotient \(Z(\sigma) / \{e, \sigma\}\) , where \(Z(\sigma)\) is the centraliser of \(\sigma\) in \(\operatorname{Aut}(\mathbf{A})\) .
\item
  If A is algebraic over its prime field Q then any two involutions of A are conjugate in Aut(A) and the centraliser of an involution \(\sigma\) is \(\{e, \sigma\}\) ; the group Aut(A) and the set of involutions of A have the cardinality of the continuum, and the centre of Aut(A) is \(\{e\}\) .
\item
  If A has finite transcendence degree over Q, show that the centraliser \(Z(\sigma)\) of any involution \(\sigma\) of A is countable (observe that the fixed subfield E of \(\sigma\) is countable, and any automorphism of E which fixes the elements of some transcendence basis for E over Q is the identity). If deg. \(tr_{Q}A = n\) , show that there exist at least \(n + 1\) conjugacy classes of involutions of A (if \(E \subset A\) is maximal orderable with \([A : E] = 2\) , consider the transcendence degree over Q of the subfield of E consisting of elements fixed by Aut(E), and use Ex. 24).
\item
  If A has infinite transcendence degree m over Q, show that there exist involutions \(\sigma\) of A such that \(\operatorname{Card}(Z(\sigma)) = \operatorname{Card}(\operatorname{Aut}(A)) = 2^{m}\) (use Ex. 24 to show that there exist maximal orderable fields E such that deg. \(tr_{Q}E = m\) , such that for any subset M of any transcendence basis B for E over Q, there exists an automorphism \(u_{M}\) of E such that \(u_{M}(x) \neq x\) for \(x \in M\) and \(u_{M}(x) = x\) for \(x \in B - M\) ).
\item
  If \(\deg \cdot tr_{Q}A = 1\) , show that there is only one conjugacy class of involutions of A whose centralisers are infinite. (If E is a maximal orderable extension of Q of transcendence degree 1 such that \(\operatorname{Aut}(E)\) is nontrivial, use Ex. 24 to show that E cannot be comparable to Q; if \(t \in E\) is infinitely large with respect to Q, show that so is \(u(t)\) for any automorphism u of E, and that conversely if \(t'\) is infinitely large with respect to Q then there exists an automorphism u of E such that \(u(t) = t'\) , using Ex. 15 of VI, p.~40.)
\item
  Let u be an automorphism of A which commutes with all the involutions of A. Show that if \(a \in A\) is transcendental over Q then a and \(u(a)\) are algebraically dependent over Q. (Consider two elements b, c of A such that \(b^{2} = a\) , \(c^{2} = -u(a)\) and observe that there exists a maximal orderable extension \(E \subset A\) such that \([A : E] = 2\) and b, c are in E; obtain a contradiction by noticing that \(u(E) = E\) .)
\end{enumerate}

\(g\) ) Deduce from \(f\) ) that \(u(a) = a\) for every element \(a \in \mathbf{A}\) transcendental over \(\mathbf{Q}\) (consider a maximal ordered extension E of \(\mathbf{Q}(a)\) contained in A, algebraic over \(\mathbf{Q}(a)\) and comparable to \(\mathbf{Q}\) , using Ex. 24). Conclude that \(u\) is necessarily the identity (observe that if \(a\) is algebraic over \(\mathbf{Q}\) and \(b\) is transcendental over \(\mathbf{Q}\) then \(ab\) is transcendental over \(\mathbf{Q}\) ).

\begin{enumerate}
\def\labelenumi{\arabic{enumi})}
\setcounter{enumi}{32}
\tightlist
\item
  Let \(\mathbf{K}\) be a commutative field (of arbitrary characteristic \(q\) ), and let \(l\) be a prime number. Let \(\mathbf{K}'\) be an algebraic extension of \(\mathbf{K}\) with the following two properties:
\end{enumerate}

\begin{enumerate}
\def\labelenumi{(\roman{enumi})}
\item
  every polynomial in \(\mathbf{K}[\mathbf{X}]\) whose degree is not divisible by \(l\) has a root in \(\mathbf{K}'\) ;
\item
  every polynomial in \(\mathbf{K}'[\mathbf{X}]\) of degree equal to \(l\) has a root in \(\mathbf{K}'\) .
\end{enumerate}

\begin{enumerate}
\def\labelenumi{\alph{enumi})}
\item
  Show that the perfect closure \(\mathbf{K}_1\) of \(\mathbf{K}\) is contained in \(\mathbf{K}'\) and satisfies the condition analogous to (i).
\item
  Every algebraic extension of K whose degree is not divisible by l is isomorphic to a subextension of \(K'\) (use the primitive element theorem).
\item
  Let L be a subfield of \(K'\) containing K, and let M be a Galois extension of L, whose degree is a power of l. Show that M is isomorphic to a sub-extension of \(K'\) (use the fact that an l-group is nilpotent).
\end{enumerate}

\(d\) ) Deduce that \(\mathbf{K}'\) is an algebraic closure of \(\mathbf{K}\) (if \(\mathbf{F}\) is a (separable) extension of \(\mathbf{K}\) of finite degree, consider a Galois extension \(\mathrm{F}_1\supset \mathrm{F}\) of finite degree, and the fixed subfield of \(\mathrm{F}_1\) under a Sylow \(l\) -subgroup of the Galois group of \(\mathrm{F}_1\) over \(\mathbf{K}\) ).

\begin{enumerate}
\def\labelenumi{\arabic{enumi})}
\setcounter{enumi}{33}
\tightlist
\item
  Let E be a totally ordered set. A (Dedekind) cut is a partition (S, D) of E into two nonempty subsets such that D is of the form M(A), the set of upper bounds for some nonempty subset A of E; a necessary and sufficient condition for D to have a least element x is that \(D = \{x, \rightarrow\{, \text{so } S = \} \leftarrow, x\{.\)
\end{enumerate}

In a totally ordered abelian group G, a cut (S, D) is called proper if for all e \textgreater{} 0 in G there exist \(x' \in S\) and \(x'' \in D\) such that \(x'' - x' < e\) . Show that (S, D) is proper if and only if D is invertible in the monoid \(\mathfrak{M}(G)\) of major subsets of G (VI, p.~35, Ex. 30); all cuts are proper if and only if G is archimedean (VI, p.~35, Ex. 31); and hence * isomorphic to a subgroup of R (VI, p.~36, Ex. 33). *

\begin{enumerate}
\def\labelenumi{\arabic{enumi})}
\setcounter{enumi}{34}
\tightlist
\item
  A subfield K of an ordered field E is dense in E if for every nonempty open interval \(a, b \in E\) there exists \(x \in K\) such that a \textless{} x \textless{} b. Show that E is then comparable to K (Ex. 11, a)). If E is taken to be R(X), where R is a maximal ordered field and X is infinitely large with respect to R (VI, p.~43, Ex. 24), show that E is comparable to the subfield \(K = R(X^{2})\) , but that K is not dense in E.
\end{enumerate}

For K to be dense in E, it is necessary and sufficient that for all \(x \in E\) the partition of K consisting of \(K \cap J \leftarrow, x\) and \(K \cap [x, \rightarrow]\) be a proper cut (Ex. 34).

\begin{enumerate}
\def\labelenumi{\arabic{enumi})}
\setcounter{enumi}{35}
\item
  Let K be an ordered field and let \(\tilde{K}\) be the set of major sets D such that the cut (S, D) is proper; show that \(\tilde{K}\) is an ordered field, taking addition to be that of the monoid \(\mathfrak{M}(K)\) , and the product of two elements \(D_{1}\) and \(D_{2}\) of \(\tilde{K}\) contained in \(K_{+}\) to be the set \(\langle D_{1}D_{2}\rangle\) (VI, p.~35, Ex. 30), where \(D_{1}D_{2}\) denotes the set of \(x_{1}x_{2}\) for \(x_{1}\in D_{1}\) and \(x_{2}\in D_{2}\) . Show that for any increasing homomorphism f from K onto a dense subfield F of an ordered field E, there exists an increasing homomorphism g from E into \(\tilde{K}\) such that \(g\circ f\) is the natural injection of K into \(\tilde{K}\) .
\item
  Let K be an ordered field and let n \textgreater{} 0 be an integer such that every element x \textgreater{} 0 of K is equal to a power \(y^{n}\) of some element y \textgreater{} 0. Show that the field \(\tilde{K}\) has the same property.
\item
  Suppose that the field K is maximal ordered. Show that the field \(\tilde{\mathbf{K}}\) is also maximal ordered. (One can argue by contradiction, by considering an ordered algebraic extension E of \(\tilde{\mathbf{K}}\) which is maximal ordered, and supposing that \(\mathrm{E} \neq \tilde{\mathbf{K}}\) . Then there exists among the elements of \(\mathrm{E} \cap [\tilde{\mathbf{K}}\) an element \(w\) whose minimal polynomial \(f \in \tilde{\mathbf{K}}[\mathbf{X}]\) has least degree \(n > 1\) . Show that there exists an interval \(\mathbf{a}, b\) (in E such that \(a, b \in \mathbf{K}\) and \(a < w < b\) , and an element \(e > 0\) of E such that \(f'(x) > e\) for all \(x \in [\mathbf{a}, b]\) , or \(f'(x) < -e\) for all \(x \in [\mathbf{a}, b]\) . Show that there exists \(r > 0\) in K such that \(f(a) \leqslant -r\) and \(r \leqslant f(b)\) , or \(f(a) \geqslant r\) and \(f(b) \leqslant -r\) . Approximate the coefficients of \(f\) by elements of K, and show that we obtain in this way a polynomial \(g \in \mathbf{K}[\mathbf{X}]\) such that the function \(x \mapsto g(x)\) is monotone in \([\mathbf{a}, \mathbf{b}]\) and vanishes at some point of K in this interval, arbitrarily close to \(w\) .)
\item
  Let K be an ordered field and let E be an ordered algebraic extension which is maximal ordered. Show that the only K-automorphism of the field E is the identity. (First show that every K-homomorphism from E to itself is surjective, and is an isomorphism of ordered fields; if \(x \in E\) has minimal polynomial \(f \in K[X]\) then an automorphism of E necessarily permutes the roots of f belonging to E. Deduce that it fixes x.)
\item
  Let K be an ordered field. Make the field of formal power series \(E = K((X))\) an ordered field by taking the elements \textgreater0 to be the formal power series \(\sum_{n=-h}^{\infty} r_n X^n\) whose
\end{enumerate}

nonzero coefficient of least degree is \textgreater0.

\begin{enumerate}
\def\labelenumi{\alph{enumi})}
\item
  Show that if \(L = K(X)\) is the subfield of E consisting of rational functions in one indeterminate, then \(E = \tilde{L}\) (Ex. 36) up to isomorphism of ordered fields (observe that for every element a \textgreater{} 0 of L there exists an integer n \textgreater{} 0 such that \(0 < X^{n} < a\) ).
\item
  Suppose that K is maximal ordered. Show that the union F of the fields \(\mathbf{K}((\mathbf{X}^{1/n}))\) (ordered in the same way as E) is an ordered algebraic extension of E which is maximal ordered (use Ex. 2 of V, p.~150).
\item
  Show that the field \(\widetilde{\mathbf{F}}\) (Ex. 36) is the field of formal series \(\sum_{r} \alpha_r X^r\) , where \(r\) runs through
\end{enumerate}

the set Q of rational numbers, and where \(\alpha_{r} \in K\) and the set of r such that \(\alpha_{r} \neq 0\) is an increasing sequence which is either finite or tends to \(+\infty\) ; order this field in the same manner as E. In particular \(\widetilde{F} \neq F\) .

\begin{enumerate}
\def\labelenumi{\arabic{enumi})}
\setcounter{enumi}{40}
\tightlist
\item
  Let K be an ordered field, and E a vector space of dimension 2 over K.
\end{enumerate}

\begin{enumerate}
\def\labelenumi{\alph{enumi})}
\item
  Show that the set of triples of pairwise distinct lines ( \(D_{1}\) , \(D_{2}\) , \(D_{3}\) ) passing through 0 has two orbits under the action of the group \(GL^{+}(E)\) (consider the subgroup fixing two lines).
\item
  The set of triples of pairwise distinct half-lines ( \(\Delta_{1}\) , \(\Delta_{2}\) , \(\Delta_{3}\) ) with origin 0 has 7 orbits under the action of \(GL(E)\) , and 14 under the action of \(GL^{+}(E)\) (same method).
\end{enumerate}

